\vspace{-3mm}
% --- [科研经历 / Research Experience] ---
\cvsection{\textcolor{awesome}{科研}经历}
\vspace{-3mm}
\begin{cventries}
  \cventry
    {研究实习生 · 自回归视频生成方向} % 角色/方向
    {北京大学 IfLab 课题组} % 机构
    {远程 / 北京} % 地点【待确认】
    {2025.0X - 至今} % 时间【待确认】
    {
      \vspace{-4mm}
      \begin{cvitems}
        \item {\textbf{Pyramid Forcing（主导课题，预投 NeurIPS）}——面向 Video World Model 长时序生成的 Head-level 动态 KV Cache 优化，独立负责 Idea 提出、Method 实现与 Experiments 评测全流程。}
        \item {\textbf{[背景]} 针对 Self-Forcing 类长序列自回归视频生成在长帧率下的【待补充：注意力退化 / 显存墙】现象——该问题在长时序推理中直接构成可用性瓶颈。}
        \item {\textbf{[洞察]} 从第一性原理剖析其 Attention 循环机制，【待补充：填入你真正推导出的数学结论；能白板推导才写「推导并证实其与 RoPE 的交互性质」，否则降级为「实验定位到周期性衰减」】。}
        \item {\textbf{[行动]} 据此在代码层设计【待补充：Head-level 动态 KV Cache 级联淘汰策略】并重写注意力缓存调度。}
        \item {\textbf{[成果]} 在【待补充：FVD/SSIM】视觉指标衰减低于【待补充：0.X\%】的约束下，将【待补充：1024】帧推理的 \textbf{KV Cache 显存占用降低【待补充：XX\%】}。}
        \item {\textbf{t2v-eval（评测平台）}——搭建 \textbf{Vbench-Long} 评测环境与实验结果可视化流水线，支撑 Pyramid Forcing 的指标对比与消融分析。}
        \item {\textbf{Stanford CS336 · Assignment 1（入组考核）}——不依赖高层封装，从零实现 Transformer 训练全链路（tokenizer / 架构 / 训练循环），训练可生成 120-150 词连贯文本的 GPT。}
      \end{cvitems}
    }
\end{cventries}
